% Solutions for chapters/14b-semiclassical-index.tex
\providecommand{\Kc}{K_{\mathrm{c}}}
\providecommand{\inda}{\operatorname{ind}_{\mathrm{a}}}
\providecommand{\Ginf}{G^{-\infty}}
\providecommand{\Indsl}[1]{\operatorname{Ind}^{#1}_{\mathrm{sl}}}
\providecommand{\indT}{\operatorname{ind}_{\mathcal T}}
\providecommand{\Psisl}[1]{\Psi^{#1}_{\mathrm{sl}}}
\providecommand{\Psiinf}{\Psi^{-\infty}_{\epsilon^\infty}}
\providecommand{\sigsl}{\sigma_{\mathrm{sl}}}
\providecommand{\Ope}{\operatorname{Op}_\epsilon}
\providecommand{\Torus}{\mathbb{T}}

\section{Solutions for \cref{ch:semiclassical-index}: The index theorem via semiclassical and adiabatic limits}\label{solutions:semiclassical-index}

Throughout, for a compactly supported idempotent perturbation $e$ of $p_0$ we write
$g_{e,p_0}$ for the loop of \cref{scindex:def:d}. Several solutions use the
following computation of the relative index when the base is a point; we prove it
first.

\begin{lemma}\label{solsemiclassical-index:lem:d-trace}
  Let $B=\mathrm{pt}$, so that $M$ is a single (possibly non-compact) manifold,
  and let $e=p_0+R$ be a compactly supported idempotent perturbation of $p_0$
  (\cref{scindex:def:d}). Then $R$ is of trace class on $L^2(M;E)$ and, under
  $K^0(\mathrm{pt})=\Z$,
  \begin{equation}\label{solsemiclassical-index:eq:d-trace}
    d(e,p_0)=\Tr(e-p_0).
  \end{equation}
\end{lemma}
\begin{proof}
  \emph{Trace class.} $M$ is an open subset of a compact manifold $\bar M$, and
  the kernel of $R$ lies in $\CIc(M\times M;\Hom(E)\otimes\Omega_R)$; extending
  $E$ over a neighbourhood of the support and the kernel by $0$, $R$ becomes a
  smoothing operator on a compact manifold, hence of trace class (as used in the
  proof of \cref{eta:lem:trace-class}).

  \emph{The Toeplitz index on a point is minus the winding number.} By
  \ref{scindex:K-Ginf} and \cref{bott:lem:stable}, $K^{-2}(\mathrm{pt})$ is the
  group of stable pointed homotopy classes of loops in the groups $\GL(N,\C)$. We
  use the standard fact that $\pi_1(\GL(N,\C))\cong\Z$, the isomorphism being the
  winding number of the determinant. The winding number of \ref{bott:B:wind},
  $\wind(h)=\frac1{2\pi i}\int_0^{2\pi}\Tr(h^{-1}\pa_\theta h)\,d\theta$, is
  homotopy invariant, is additive under products of loops, since
  $\Tr((h_1h_2)^{-1}\pa_\theta(h_1h_2))=\Tr(h_1^{-1}\pa_\theta h_1)+
  \Tr(h_2^{-1}\pa_\theta h_2)$ by cyclicity of the trace, and equals the winding
  number of $\det h$ on loops in $\GL(V)$. Hence $\wind:K^{-2}(\mathrm{pt})\to\Z$ is
  an isomorphism, and $K^{-2}(\mathrm{pt})$ is generated by the class of
  $\ell=pz^{-1}+\Id-p$, $p$ a constant projection of rank one, for which
  $\det\ell=z^{-1}$ and $\wind(\ell)=-1$. By \eqref{scindex:eq:bott-map},
  $\indT[\ell]=[\Ran p]=1$. Two homomorphisms $\Z\to\Z$ which agree on a generator
  are equal, so $\indT=-\wind$ on $K^{-2}(\mathrm{pt})$.

  \emph{The winding number of $g_{e,p_0}$.} The class of $g=g_{e,p_0}$ is that of
  $\rho_J(g)=\Id+JAJ^*$, $g=\Id+A$, for an admissible embedding $J$
  (\cref{scindex:def:class}). Since $J^*J=\Id$, we have $\rho_J(g)^{-1}=\Id+JA'J^*$
  with $g^{-1}=\Id+A'$, and
  \begin{equation*}
    \rho_J(g)^{-1}\pa_\theta\rho_J(g)=(\Id+JA'J^*)J(\pa_\theta A)J^*=
    J(\Id+A')(\pa_\theta A)J^*=J\bigl(g^{-1}\pa_\theta g\bigr)J^*,
  \end{equation*}
  whose trace is $\Tr(g^{-1}\pa_\theta g\,J^*J)=\Tr(g^{-1}\pa_\theta g)$. Now
  $g=ac$ with $a=ez^{-1}+\Id-e$, $c=p_0z+\Id-p_0$, $z=e^{i\theta}$, and
  $a^{-1}=ez+\Id-e$, $c^{-1}=p_0z^{-1}+\Id-p_0$. Hence
  $a^{-1}\pa_\theta a=(ez+\Id-e)(-iz^{-1}e)=-ie$ and $c^{-1}\pa_\theta c=
  (p_0z^{-1}+\Id-p_0)(izp_0)=ip_0$, so
  \begin{equation*}
    g^{-1}\pa_\theta g=c^{-1}(a^{-1}\pa_\theta a)c+c^{-1}\pa_\theta c
    =-ic^{-1}(e-p_0)c,
  \end{equation*}
  using that $c$ commutes with $p_0$. This is of trace class with trace
  $-i\Tr(e-p_0)$. Therefore $\wind(g)=\frac1{2\pi i}\cdot2\pi\cdot(-i)\Tr(e-p_0)
  =-\Tr(e-p_0)$, and $d(e,p_0)=\indT[g]=-\wind(g)=\Tr(e-p_0)$.
\end{proof}

%----------------------------------------------------------------------
\subsection*{\Cref{scindex:ex:strong}}

\emph{Idea:} $U_t$ is $U$ squeezed into the interval $[0,t]$; as $t\to0$ it acts
on less and less of the function, and the squeezing is continuous in the strong
topology.

\emph{Unitarity.} For $v\in L^2([0,t])$,
$\norm{\kappa_tv}^2=\int_0^1t\abs{v(tx)}^2dx=\int_0^t\abs{v(y)}^2dy$, and
$\kappa_t$ is onto, with inverse $(\kappa_t^{-1}w)(y)=t^{-1/2}w(y/t)$. So
$\kappa_t$ is unitary, $\kappa_t^{-1}U\kappa_t$ is a unitary of $L^2([0,t])$, and
$U_t$ is unitary on $H=L^2([0,t])\oplus L^2([t,1])$ as a direct sum of unitaries.
$U_0=\Id$ by definition.

\emph{The estimate.} Write $v=v_1+v_2$, $v_1=1_{[0,t]}v$. Then $U_tv-v=
\kappa_t^{-1}U\kappa_tv_1-v_1$, of norm at most $2\norm{v_1}=2\norm v_{L^2([0,t])}$.

\emph{Multiplicativity.} For $t>0$ both $U_tV_t$ and $(UV)_t$ are block diagonal,
with first blocks $\kappa_t^{-1}U\kappa_t\kappa_t^{-1}V\kappa_t=\kappa_t^{-1}UV\kappa_t$
and second blocks $\Id$; for $t=0$ both sides are $\Id$.

\emph{Continuity.} For $t\in(0,1]$ define $S_t:H\to H$ by $(S_tw)(x)=t^{1/2}w(tx)$.
Its adjoint is $(S_t^*u)(y)=t^{-1/2}u(y/t)$ for $y<t$ and $0$ for $y\ge t$, since
$\int_0^1t^{1/2}w(tx)\bar u(x)\,dx=\int_0^tw(y)\,t^{-1/2}\bar u(y/t)\,dy$. Thus
$S_t^*S_t$ is multiplication by $1_{[0,t]}$, $\norm{S_t}=\norm{S_t^*}=1$, and
\begin{equation*}
  U_t=S_t^*US_t+\bigl(\Id-1_{[0,t]}\bigr),\qquad t\in(0,1].
\end{equation*}
The maps $t\mapsto S_tw$ and $t\mapsto S_t^*u$ are continuous on $(0,1]$ for each
$w,u\in H$. Indeed, for $w$ continuous on $[0,1]$, $S_tw\to S_{t_0}w$ uniformly as
$t\to t_0>0$, by uniform continuity of $w$; for $u$ continuous on $[0,1]$ with
$u(1)=0$, extend $u$ by $0$ to a uniformly continuous function $\tilde u$ on
$[0,\infty)$, so that $S_t^*u(y)=t^{-1/2}\tilde u(y/t)\to S_{t_0}^*u(y)$
uniformly. These functions are dense in $H$ and the operators have norm $1$, so a
$\frac\delta3$-argument gives strong continuity. Also $t\mapsto1_{[0,t]}v$ is
continuous, since $\norm{1_{[0,t]}v-1_{[0,t_0]}v}^2$ is the integral of $\abs v^2$
between $t$ and $t_0$.

Fix $(U_0,t_0,v_0)$. For all $(U,t,v)$, since $U_t$ is unitary,
\begin{equation}\label{solsemiclassical-index:eq:cont1}
  \norm{U_tv-(U_0)_{t_0}v_0}\le\norm{v-v_0}+\norm{U_tv_0-(U_0)_{t_0}v_0}.
\end{equation}
If $t_0=0$, the last term is $\norm{U_tv_0-v_0}\le2\norm{v_0}_{L^2([0,t])}$, which
tends to $0$ as $t\to0$ uniformly in $U$. If $t_0>0$, then for $t>0$, with
$w_0=S_{t_0}v_0$ and $u_0=U_0w_0$,
\begin{multline*}
  U_tv_0-(U_0)_{t_0}v_0=S_t^*U(S_tv_0-S_{t_0}v_0)+S_t^*(U-U_0)w_0
  +(S_t^*-S_{t_0}^*)u_0\\+(1_{[0,t_0]}-1_{[0,t]})v_0,
\end{multline*}
whose norm is at most $\norm{S_tv_0-S_{t_0}v_0}+\norm{(U-U_0)w_0}+
\norm{(S_t^*-S_{t_0}^*)u_0}+\norm{(1_{[0,t_0]}-1_{[0,t]})v_0}$. The second term is
small on the strong neighbourhood $\{U;\ \norm{(U-U_0)w_0}<\delta\}$ of $U_0$, the
others for $t$ near $t_0$. With \eqref{solsemiclassical-index:eq:cont1} this proves
joint continuity of $(U,t,v)\mapsto U_tv$.

\emph{Contractibility.} Define $F:\mathrm U(H)\times[0,1]\to\mathrm U(H)$ by
$F(U,t)=U_t$. The strong topology is the coarsest topology making all maps
$U\mapsto Uv$ continuous, so $F$ is continuous into $\mathrm U(H)$ with the strong
topology if and only if each $(U,t)\mapsto U_tv$ is continuous, which we have
shown. Since $\kappa_1=\Id$, $F(U,1)=U$, and $F(U,0)=\Id$. So $F$ is a homotopy
from the identity of $\mathrm U(H)$ to a constant map: $\mathrm U(H)$ is
contractible in the strong topology. (Since $(UV)_t=U_tV_t$, the contraction is even
through group homomorphisms.) Any separable infinite dimensional Hilbert space is
unitarily equivalent to $L^2([0,1])$, and conjugation by a fixed unitary is a
homeomorphism for the strong topologies, so the same holds for its unitary group.

%----------------------------------------------------------------------
\subsection*{\Cref{scindex:ex:frame}}

\emph{Idea:} finitely many smooth sections whose restrictions to each fibre are
linearly independent span a smooth vector bundle, and its orthogonal projection is
a smooth family of finite rank smoothing operators. Finite rank approximation on
the torus gives sections which approximately span everything a smoothing operator
sees; a small generic perturbation makes them independent on every fibre.

If $\dim Z=0$ the fibres are finite sets, $L^2(Z_b;E_b)$ is finite dimensional,
$\Id\in\Psido{-\infty}(M/B;E)$, and $\pi_k=\Id$ does everything. Assume $\dim Z\ge1$.
We may also assume $B$ connected (otherwise treat each component).

\emph{A remark on the hint of the original statement.} One could try to take for
$\pi_k$ the Riesz projection near $1$ of the finite rank self-adjoint family
$T_b=J_b^*P_VJ_b$. This requires a circle around $1$ which meets the spectrum of no
$T_b$, and such a gap is not available in general: the non-zero eigenvalues of
$T_b$ are those of the Hermitian matrix $P_VJ_bJ_b^*P_V$ on $V$, which vary
continuously with $b$ and may cross any given level as the embedded fibres move. The
statement has been amended accordingly; the construction below uses the operators
$T_b$ only to produce sections.

\emph{(a) Step 1: sections from the embedding.} Let $J$ be the admissible embedding
of \cref{scindex:lem:embedding}, with $K=W=M$ ($M$ is compact). In its
construction $Y=\Torus^N$, $p_b:\Omega_b\to Z_b$ is the projection of a tubular
neighbourhood, and $dy=\mu_{b,m}\otimes\nu(m)$ on $\Omega_b$ for a family of smooth
positive measures $\mu_{b,m}$ on the fibres $p_b^{-1}(m)$, depending smoothly on
$(b,m)$. Hence, from $J_bu(y)=\gamma(y)\kappa(u(p_b(y)))$,
\begin{equation*}
  (J_b^*f)(m)=\kappa(m)^*\int_{p_b^{-1}(m)}\gamma(y)f(y)\,d\mu_{b,m}(y),
  \qquad f\in L^2(\Torus^N;\C^R),
\end{equation*}
and for $f\in\CI(\Torus^N;\C^R)$ the function $m\mapsto(J^*_{\phi(m)}f)(m)$ is a
smooth section of $E$ over $M$. Let $V_1\subset V_2\subset\dots$ be the subspaces of
\ref{bott:B:finite} on $(\Torus^N,\C^R)$, with orthogonal projections $P_k$, and
let $e_1,\dots,e_{n_k}$ be an orthonormal basis of $V_k$. Put
$t_j=J^*e_j\in\CI(M;E)$, so that
\begin{equation*}
  T_ku=\sum_{j\le n_k}\inner u{t_j}t_j=J^*P_kJu.
\end{equation*}
For $A\in\Psido{-\infty}(M/B;E)$ the family $X_b=J_bA_bJ_b^*$ is a smooth map
$B\to\Psido{-\infty}(\Torus^N;\C^R)$, so by \ref{bott:B:finite}
$\delta_k=\sup_b\norm{P_kX_bP_k-X_b}\to0$. Since $(\Id-P_k)P_k=0$,
$(\Id-P_k)X=(\Id-P_k)(X-P_kXP_k)$ and $X(\Id-P_k)=(X-P_kXP_k)(\Id-P_k)$, so both
have norm at most $\delta_k$. As $J^*J=\Id$ we have $A=J^*XJ$ and
$(\Id-T_k)A=J^*(\Id-P_k)XJ$, $A(\Id-T_k)=J^*X(\Id-P_k)J$; hence
\begin{equation}\label{solsemiclassical-index:eq:Tk}
  \sup_b\norm{(\Id-T_k)A}\le\delta_k,\qquad\sup_b\norm{A(\Id-T_k)}\le\delta_k .
\end{equation}

\emph{Step 2: a genericity lemma.} \textit{Given $s_1,\dots,s_n\in\CI(M;E)$ and
$\eta>0$, there are $s'_j\in\CI(M;E)$ with $\sup_M\abs{s'_j-s_j}<\eta$ whose
restrictions to every fibre $Z_b$ are linearly independent.}

Let $r=\operatorname{rank}E$, $d_B=\dim B$, and choose $K\in\N$ with
$2rK>d_B+2n$ and distinct points $z_1,\dots,z_K$ of the typical fibre $Z$ (possible
since $\dim Z\ge1$). Cover $B$ by finitely many open sets $U_i$, each
diffeomorphic to a ball or half-ball and carrying a trivialization
$f_i:\phi^{-1}(U_i)\to Z\times U_i$ (\cref{product:thm:local-triviality}), and choose
compact $K_i\subset U_i$ covering $B$. Put $m_{i,k}(b)=f_i^{-1}(z_k,b)$; for
fixed $i$ and $b$ these are $K$ distinct points of $Z_b$. The pull-backs of $E$ by
the maps $m_{i,k}:U_i\to M$ are trivial ($U_i$ is contractible), so evaluation
gives linear maps $\mathrm{ev}_i(b):\CI(M;E)\to\C^{rK}$, $s\mapsto(s(m_{i,k}(b)))_k$,
depending smoothly on $b$. Since the points are distinct, for each $b_0\in K_i$
there are $rK$ sections, supported near the points $m_{i,k}(b_0)$, whose values
$\mathrm{ev}_i(b_0)$ form a basis of $\C^{rK}$; by continuity they still do for $b$
near $b_0$. By compactness there are finitely many sections $w_1,\dots,w_L$ and
open sets $U'_i$, $K_i\subset U'_i\subset U_i$, such that
$\{\mathrm{ev}_i(b)w_l\}_l$ spans $\C^{rK}$ for all $i$ and $b\in U'_i$.

For $c=(c_{jl})\in\C^{n\times L}$ put $s_j^c=s_j+\sum_lc_{jl}w_l$ and consider
\begin{equation*}
  F_i:U'_i\times(\C^n\setminus0)\times\C^{n\times L}\to\C^{rK},\qquad
  F_i(b,\lambda,c)=\mathrm{ev}_i(b)\Bigl(\sum_j\lambda_js^c_j\Bigr).
\end{equation*}
Its derivative in $c$ alone, $\delta c\mapsto\sum_l(\sum_j\lambda_j\delta c_{jl})
\mathrm{ev}_i(b)w_l$, is onto because $\lambda\ne0$ (take $\delta c_{jl}=
\bar\lambda_j\mu_l/\abs\lambda^2$). So $F_i$ restricted to
$(U'_i\setminus\pa B)\times(\C^n\setminus0)\times\C^{n\times L}$, and to
$(U'_i\cap\pa B)\times(\C^n\setminus0)\times\C^{n\times L}$, is a submersion of a
manifold without boundary, and the zero sets of these restrictions are submanifolds
of dimension at most $d_B+2n+2nL-2rK<2nL$. The image of a manifold of dimension
less than $2nL$ under a smooth map to $\C^{n\times L}\cong\R^{2nL}$ has measure zero
(Sard's theorem). Hence there is a null set $\mathcal N\subset\C^{n\times L}$ such
that for $c\notin\mathcal N$, $F_i(b,\lambda,c)\ne0$ for all $i$, $b\in U'_i$,
$\lambda\ne0$. Choose such $c$ so small that $\sup_M\abs{s^c_j-s_j}<\eta$, and put
$s'_j=s_j^c$. If $\sum_j\lambda_js'_j$ vanished on $Z_b$ for some $\lambda\ne0$,
then, with $b\in K_i$, it would vanish at the points $m_{i,k}(b)$, i.e.\
$F_i(b,\lambda,c)=0$, which is excluded.

\emph{Step 3: the projections.} Apply Step 2 to $t_1,\dots,t_{n_k}$, with $\eta$ so
small that $T'_ku=\sum_j\inner u{s'_j}s'_j$ satisfies $\sup_b\norm{T'_k-T_k}\le1/k$
(the $L^2(Z_b)$ norms are bounded by a constant times the sup norms, uniformly in
$b$, and $\norm{T'_k-T_k}\le\sum_j(\norm{s'_j-t_j}\norm{s'_j}+\norm{t_j}
\norm{s'_j-t_j})$). Let $S_b:\C^{n_k}\to L^2(Z_b;E_b)$, $S_bc=\sum_jc_js'_j|_{Z_b}$.
It is injective for every $b$, so the Gram matrix $G_b=S_b^*S_b$, with entries
$\inner{s'_j}{s'_i}_{L^2(Z_b)}$, is a smooth family of invertible matrices, and
\begin{equation*}
  \pi_k=S(S^*S)^{-1}S^*,\qquad
  \pi_k(m,m')=\sum_{i,j}s'_i(m)\,(G_b^{-1})_{ij}\,s'_j(m')^*\,\nu(m'),
\end{equation*}
is a smooth family of orthogonal projections of rank $n_k$ onto
$N_b=\operatorname{span}\{s'_j|_{Z_b}\}$, with smooth kernel: $\pi_k\in
\Psido{-\infty}(M/B;E)$. Since $T'_k=SS^*$ is self-adjoint with range $N_b$,
$(\Id-\pi_k)T'_k=0=T'_k(\Id-\pi_k)$. Therefore, using
\eqref{solsemiclassical-index:eq:Tk},
\begin{equation*}
  \norm{(\Id-\pi_k)A}=\norm{(\Id-\pi_k)(A-T'_kA)}\le\delta_k+\norm A/k,
\end{equation*}
and likewise $\norm{A(\Id-\pi_k)}\le\delta_k+\norm A/k$, uniformly in $b$
($\sup_b\norm{A_b}<\infty$ by \cref{scindex:lem:uniform}, or directly). Finally
$\pi_kA\pi_k-A=-(\Id-\pi_k)A-\pi_kA(\Id-\pi_k)$, so
$\sup_b\norm{\pi_kA\pi_k-A}\le2(\delta_k+\norm A/k)\to0$.

\emph{(b) Melrose's construction.} Cover $B$ by open sets $U_i$ with trivializations
$f_i$ of the fibration and isomorphisms of $E$ with the pull-back of a bundle $E_0$
over $Z$ (possible on contractible $U_i$). They give unitaries
$W_{i,b}:L^2(Z_b;E_b)\to H_0=L^2(Z;E_0)$, $b\in U_i$ (pull-back by
$f_i(\cdot,b)^{-1}$, times the square root of the ratio of the densities), which map
smooth sections of $E$ over $\phi^{-1}(U_i)$ to smooth families in $\CI(Z;E_0)$
and back, and conjugate smooth families of smoothing operators to smooth families
in $\Psido{-\infty}(Z;E_0)$. The transition maps $g_{ij}(b)=W_{i,b}W_{j,b}^{-1}$
are unitaries of $H_0$ induced by diffeomorphisms of $Z$; they depend continuously
on $b$ in the strong topology (on smooth sections they are continuous in the
$\CI$ topology, and they are uniformly bounded). The bundle $\mathcal F$ whose fibre
at $b$ is the set of unitaries $H_0\to L^2(Z_b;E_b)$, trivialized over $U_i$ by
$U\mapsto W_{i,b}U$, with transition maps given by left multiplication by
$g_{ij}$ (continuous for the strong topology, because multiplication is jointly
strongly continuous on the unitary group), is a fibre bundle with fibre
$\mathrm U(H_0)$. Since $\dim Z\ge1$, $H_0$ is separable and infinite dimensional,
so by \cref{scindex:ex:strong} this fibre is contractible, and by the fact quoted
in the exercise $\mathcal F$ has a continuous section $b\mapsto U_b$. Let
$\{f_i\}_{i\ge1}$ be an orthonormal basis of $H_0$ and $e'_i(b)=U_bf_i$: these are
continuous sections of the Hilbert bundle (in each trivialization
$b\mapsto W_{i,b}e'_i(b)$ is norm continuous) forming an orthonormal basis of every
fibre.

\emph{Smoothing the sections.} Let $e'$ be a continuous section and $\delta>0$. For
$b_0\in U_i$ choose $\varphi\in\CI(Z;E_0)$ with
$\norm{\varphi-W_{i,b_0}e'(b_0)}<\delta$. Then $\sigma(b)=W_{i,b}^{-1}\varphi$ is a
smooth section of $E$ over $\phi^{-1}(U_i)$ and, by continuity,
$\norm{\sigma(b)-e'(b)}=\norm{\varphi-W_{i,b}e'(b)}<\delta$ on a neighbourhood
$O_{b_0}$ of $b_0$. Cover $B$ by finitely many such $O_\alpha$, with sections
$\sigma_\alpha$, and let $\rho_\alpha$ be a subordinate partition of unity. Then
$e''=\sum_\alpha\rho_\alpha\sigma_\alpha\in\CI(M;E)$ and
$\norm{e''(b)-e'(b)}\le\sum_\alpha\rho_\alpha(b)\norm{\sigma_\alpha(b)-e'(b)}<\delta$.
Doing this with $\delta=2^{-i-4}$ for $e'_i$ gives smooth $e''_i$ with
$\sup_b\norm{e'_i-e''_i}\le2^{-i-4}$.

\emph{The operator $T_b$.} For $u\in L^2(Z_b;E_b)$,
\begin{equation*}
  \norm{(T_b-\Id)u}=\Bigl\|\sum_i\inner u{e'_i}(e''_i-e'_i)\Bigr\|
  \le\sum_{i\ge1}\abs{\inner u{e'_i}}\,2^{-i-4}\le\norm u\sum_{i\ge1}2^{-i-4}
  =\frac{\norm u}{16},
\end{equation*}
so the series converges and $\norm{T_b-\Id}\le\frac1{16}\le\frac14$. In
particular $T_b$ is invertible, and $T_be'_i=e''_i$: the $e''_i(b)$ are the image of
an orthonormal basis under an invertible operator. Hence every finite family
$e''_1(b),\dots,e''_k(b)$ is linearly independent, and their span is dense.

\emph{Gram--Schmidt.} Define $f_1,f_2,\dots$ by the Gram--Schmidt procedure applied
to $e''_1,e''_2,\dots$ fibrewise. At step $k$ one divides
$e''_k-\sum_{i<k}\inner{e''_k}{f_i}f_i$ by its norm, which is positive by linear
independence; the inner products $\inner{e''_k}{f_i}_{L^2(Z_b)}$ are fibre
integrals of smooth sections, hence smooth in $b$, so inductively each $f_k$ is a
smooth section of $E$ over $M$. The $f_i(b)$ are orthonormal and have the same span
as the $e''_i(b)$, which is dense: they form a smooth orthonormal basis of the
Hilbert bundle.

\emph{Back to (a).} With this basis, $\pi_k=\sum_{i\le k}f_i\otimes f_i^*$ is a smooth
family of finite rank orthogonal projections in $\Psido{-\infty}(M/B;E)$. For
$A\in\Psido{-\infty}(M/B;E)$ the functions $h_k(b)=\norm{(\Id-\pi_k)A_b}$ are
continuous (in a local trivialization both families are norm continuous), they
decrease in $k$ because $\pi_{k+1}\ge\pi_k$, and $h_k(b)\to0$ for each $b$ because
$\pi_k\to\Id$ strongly and $A_b$ is compact. By Dini's theorem $h_k\to0$ uniformly
on the compact set $B$; the same holds for $\norm{A_b(\Id-\pi_k)}=
\norm{(\Id-\pi_k)A_b^*}$, and as in (a), $\sup_b\norm{\pi_kA\pi_k-A}\to0$.

%----------------------------------------------------------------------
\subsection*{\Cref{scindex:ex:relative}}

\emph{Idea:} the pair $A=P-Q$, $C=\Id-P-Q$ anticommutes and satisfies
$A^2+C^2=\Id$; this pairs the eigenvalues $\pm\lambda$ of $A$ in $(-1,1)$, and only
the eigenvalues $\pm1$, which are the null space and cokernel of $QP$, are left.

\emph{Fredholm property.} Let $T=QP|_{\Ran P}:\Ran P\to\Ran Q$ and
$S=PQ|_{\Ran Q}:\Ran Q\to\Ran P$. Then
\begin{equation*}
  ST=PQP|_{\Ran P}=\Id_{\Ran P}-P(P-Q)P|_{\Ran P},\qquad
  TS=QPQ|_{\Ran Q}=\Id_{\Ran Q}+Q(P-Q)Q|_{\Ran Q},
\end{equation*}
and the correction terms are compact. By Atkinson's theorem $T$ is Fredholm.

\emph{The index.} Put $A=P-Q$ and $C=\Id-P-Q$, both self-adjoint. A direct
computation gives
\begin{equation*}
  AC=QP-PQ=-CA,\qquad A^2+C^2=(P+Q-PQ-QP)+(\Id-P-Q+PQ+QP)=\Id .
\end{equation*}
We have $\Nul T=\Ran P\cap\Nul Q=\Nul(A-\Id)$: if $Pu=u$, $Qu=0$ then $Au=u$, and
conversely $Au=u$ gives $\norm{Pu}^2-\norm{Qu}^2=\norm u^2$, so $Qu=0$ and $Pu=u$.
The orthogonal complement of $\Ran T$ in $\Ran Q$ is $\{v\in\Ran Q;\ \inner v{QPu}=0\
\forall u\in\Ran P\}=\Ran Q\cap\Nul P=\Nul(A+\Id)$, by the same argument. Hence
\begin{equation*}
  \ind T=\dim\Nul(A-\Id)-\dim\Nul(A+\Id).
\end{equation*}
If $\lambda\in(-1,1)\setminus\{0\}$ and $Au=\lambda u$, then $A(Cu)=-CAu=-\lambda Cu$
and $C^2u=(1-\lambda^2)u\ne0$ for $u\ne0$; so $C$ maps $\Nul(A-\lambda)$ injectively
into $\Nul(A+\lambda)$, and by symmetry these eigenspaces, finite dimensional since
$A$ is compact, have equal dimensions.

Now let $P$, $Q$ be as in \cref{scindex:def:d} with $B$ a point: $Q=p_0$ is an
orthogonal projection acting by multiplication, and $P=p_0+R$ is an orthogonal
projection with $R$ smoothing with compactly supported kernel. By
\cref{solsemiclassical-index:lem:d-trace}, $A=R$ is of trace class, so its trace is
the absolutely convergent sum of its eigenvalues with multiplicity (spectral theorem
for compact self-adjoint operators). The eigenvalues lie in $[-1,1]$
($\norm A\le1$ because $A^2\le A^2+C^2=\Id$), those in $(-1,1)\setminus\{0\}$ cancel
in pairs, and
\begin{equation*}
  \Tr(P-Q)=\dim\Nul(A-\Id)-\dim\Nul(A+\Id)=\ind T .
\end{equation*}
By \cref{solsemiclassical-index:lem:d-trace}, $d(P,Q)=\Tr(P-Q)$, so
$\ind(QP:\Ran P\to\Ran Q)=d(P,Q)$. (The route suggested in the exercise also works: both sides are invariant under
smooth deformations of $P$ through compactly supported perturbations of $Q$ which
are orthogonal projections, and in the model case $P=Q+\Pi$, $\Pi$ an orthogonal
projection of finite rank with $Q\Pi=\Pi Q=0$, every $u\in\Ran P$ is $u=q+w$ with
$q\in\Ran Q$, $w\in\Ran\Pi$, and $Tu=q$; so $T$ is onto with null space $\Ran\Pi$,
and $\ind T=\operatorname{rank}\Pi=d(P,Q)$ by \ref{scindex:d-finite}. The trace
computation above avoids the deformation.)

%----------------------------------------------------------------------
\subsection*{\Cref{scindex:ex:milnor-trace}}

\emph{Idea:} the trace of $e_{A,B}-p_0$ is read off from the diagonal blocks, and
the identities $R_+^2-R_+=-R_+BA$, $R_-^2-R_-=-R_-AB$ together with cyclicity
remove the squares.

We take $A$, $B$ bounded and $R_\pm$ of trace class. By \eqref{scindex:eq:milnor},
\begin{equation*}
  e_{A,B}-p_0=\begin{pmatrix}-R_-^2&(\Id+R_-)AR_+\\R_+B&R_+^2\end{pmatrix}
\end{equation*}
on $H_-\oplus H_+$. Each entry is a product of bounded operators with at least one
trace class factor, so $e_{A,B}-p_0$ is of trace class; computing the trace in an
orthonormal basis of $H_-\oplus H_+$ made of bases of the two summands, only the
diagonal blocks contribute:
\begin{equation*}
  \Tr(e_{A,B}-p_0)=\Tr(R_+^2)-\Tr(R_-^2).
\end{equation*}
Next, $R_+^2-R_+=R_+(R_+-\Id)=-R_+BA$ and $R_-^2-R_-=-R_-AB$. From the proof of
\cref{scindex:lem:milnor}, $R_+B=B-BAB=BR_-$, and $R_-=\Id-AB$ commutes with $AB$.
Hence $R_+BA=BR_-A$ and $R_-AB=ABR_-=A(BR_-)$, where $BR_-$ is of trace class. For
a trace class $X:H_-\to H_+$ and a bounded $Y:H_+\to H_-$, $\Tr(YX)=\Tr(XY)$ (that
is, $\Tr[Y,X]=0$ in the sense of the hint). Therefore
\begin{equation*}
  \bigl(\Tr R_+^2-\Tr R_+\bigr)-\bigl(\Tr R_-^2-\Tr R_-\bigr)
  =-\Tr\bigl((BR_-)A\bigr)+\Tr\bigl(A(BR_-)\bigr)=0,
\end{equation*}
i.e.\ $\Tr(R_+^2)-\Tr(R_-^2)=\Tr(R_+)-\Tr(R_-)$.

\emph{\Cref{scindex:prop:index-relative} for $B$ a point.} Now $M$ is compact,
$A\in\Psido0(M;E_+,E_-)$ is elliptic and $R_\pm$ are smoothing, hence of trace
class. By \cref{solsemiclassical-index:lem:d-trace}, the two identities above and
\cref{index:thm:trace-formula} (with $m=0$),
\begin{equation*}
  d(e_{A,B},p_0)=\Tr(e_{A,B}-p_0)=\Tr(R_+)-\Tr(R_-)=\Tr(\Id-BA)-\Tr(\Id-AB)
  =\ind(A).
\end{equation*}
For $B$ a point, $\inda[\sigma_0(A)]\in K^0(\mathrm{pt})=\Z$ is the integer
$\dim\mathcal N(A,S)-L=\ind(A)$ (\cref{product:def:index-bundle}, \ref{scindex:K-inda}),
since $\ind(A\oplus S)=\ind(A)+L$ and $A\oplus S$ is surjective. This is
\cref{scindex:prop:index-relative} for $B=\mathrm{pt}$.

%----------------------------------------------------------------------
\subsection*{\Cref{scindex:ex:oscillatory}}

\emph{Idea:} after the substitution $z'=z-\epsilon W$ the conjugated operator is
an integral of the kernel profile against the phase $e^{-iW\cdot df(z)}$, up to an
error which is smooth in $\epsilon$ up to $\epsilon=0$.

\emph{Normal form of the kernel.} By \ref{scindex:S-smoothing} (with $B$ a point and
one chart $\R^n$), the kernel of $A(\epsilon)$ is supported in $K\times K$, $K$
compact, is $O(\epsilon^\infty)$ in $\CI$ away from the diagonal, and near it equals
$\epsilon^{-n}K(\epsilon,z,(z-z')/\epsilon)$ with $K$ smooth and rapidly
decreasing with all derivatives in the last variable. Let $\psi\in\CIc(\R^n)$ equal
$1$ near $0$ with support in the region where this form holds, and $\phi\in\CIc$
equal to $1$ near $K$. Then
\begin{equation*}
  A(\epsilon)(z,z')=\epsilon^{-n}\tilde K\bigl(\epsilon,z,(z-z')/\epsilon\bigr)
  +r_\epsilon(z,z'),\qquad\tilde K(\epsilon,z,W)=\phi(z)\psi(\epsilon W)
  K(\epsilon,z,W),
\end{equation*}
where $\tilde K$ is smooth on $[0,1]\times\R^n\times\R^n$, supported in
$\supp\phi$ in $z$, with $\abs{\pa_\epsilon^j\pa_z^\alpha\pa_W^\beta\tilde K}\le
C\ang W^{-N}$ for all indices, $\tilde K(0,z,W)=K(0,z,W)$, and
$r_\epsilon\in\CIc(\R^{2n})$ has all derivatives $O(\epsilon^\infty)$. By
\ref{scindex:S-symbol}, $\sigsl(A)(z,\zeta)=\hat K(0,z,\zeta)=\int e^{-iW\cdot\zeta}
K(0,z,W)\,dW$.

\emph{The limit.} Let $v_\epsilon=e^{-if/\epsilon}A(\epsilon)(e^{if/\epsilon}u)$. The
contribution of $r_\epsilon$ is $e^{-if(z)/\epsilon}\int r_\epsilon(z,z')
e^{if(z')/\epsilon}u(z')\,dz'$; each $z$-derivative of order $\le k$ is bounded by
$C\epsilon^{-k}$ times sup norms of derivatives of $r_\epsilon$, which are
$O(\epsilon^\infty)$, and it is supported in a fixed compact set; so it tends to $0$
in $\CI$. In the main term substitute $z'=z-\epsilon W$:
\begin{equation*}
  v^0_\epsilon(z)=\int\tilde K(\epsilon,z,W)\,e^{i\Phi(\epsilon,z,W)}\,
  u(z-\epsilon W)\,dW,\qquad
  \Phi(\epsilon,z,W)=\frac{f(z-\epsilon W)-f(z)}\epsilon
  =-\int_0^1df(z-t\epsilon W)\cdot W\,dt .
\end{equation*}
The last formula shows that $\Phi$ is smooth on $[0,1]\times\R^{2n}$, with
$\Phi(0,z,W)=-df(z)\cdot W$. The integrand vanishes unless $z\in\supp\phi$ and
$z-\epsilon W\in\supp u$, and then all points $z-t\epsilon W$, $t\in[0,1]$, lie in a
fixed compact set $L$ (the convex hull of $\supp\phi\cup\supp u$). Hence, on the
support of the integrand, $\abs{\pa_z^\alpha\Phi}\le C_\alpha\abs W$ for
$\abs\alpha\ge1$, with $C_\alpha$ depending on the derivatives of $f$ on $L$ only,
and every $z$-derivative of the integrand is bounded by
$C(1+\abs W)^{\abs\alpha}\ang W^{-N}$, uniformly in $(\epsilon,z)$. By dominated
convergence, $\pa_z^\alpha v_\epsilon^0(z)$ is continuous on $[0,1]\times\R^n$
(including $\epsilon=0$), and it vanishes for $z\notin\supp\phi$; by uniform
continuity on $[0,1]\times\supp\phi$, $\pa^\alpha_zv^0_\epsilon\to\pa^\alpha_zv^0_0$
uniformly. At $\epsilon=0$,
\begin{equation*}
  v^0_0(z)=\int K(0,z,W)e^{-iW\cdot df(z)}\,dW\;u(z)=\sigsl(A)(z,df(z))\,u(z).
\end{equation*}
So $v_\epsilon\to\sigsl(A)(z,df(z))u(z)$ in $\CI(\R^n)$.

\emph{Invariance.} The formula characterizes $\sigsl(A)$ without reference to the
coordinates in which the kernel was written: every covector $\zeta\in T^*_z\R^n$ is
$df(z)$ for some real $f$ (a linear function), and $u$ with $u(z)\ne0$ is arbitrary.
Let $F:U\to U'$ be a diffeomorphism and $A'=F_*A$ the transported family, which by
the first part of \ref{scindex:S-invariance} is again semiclassical. With the
density (or unitary) factors, $(A'v)(F(z))=h(z)\bigl(A(g\cdot(v\circ F))\bigr)(z)$
with smooth $g,h$ satisfying $h\,g=1$ (for the bundle version, $g$ and $h$ are
mutually inverse bundle maps). For $f'$ real on $U'$ put $f=f'\circ F$. Then
\begin{equation*}
  \bigl(e^{-if'/\epsilon}A'(e^{if'/\epsilon}u')\bigr)(F(z))
  =h(z)\bigl(e^{-if/\epsilon}A(e^{if/\epsilon}g\,(u'\circ F))\bigr)(z),
\end{equation*}
and letting $\epsilon\to0$ on both sides,
$\sigsl(A')(F(z),df'(F(z)))u'(F(z))=h(z)\sigsl(A)(z,df(z))g(z)u'(F(z))$. Since
$df(z)=dF(z)^Tdf'(F(z))$,
\begin{equation*}
  \sigsl(A')(F(z),\zeta')=h(z)\,\sigsl(A)\bigl(z,dF(z)^T\zeta'\bigr)\,g(z)
  \quad\text{for all }\zeta'\in T^*_{F(z)}U',
\end{equation*}
which for scalar or unitary factors is the statement of \ref{scindex:S-invariance}:
the symbol transforms by the induced bundle map $T^*U\to T^*U'$ (and conjugation
by the bundle isomorphism). In particular the local definitions of $\sigsl$ in
overlapping charts agree, so $\sigsl(A)$ is a well-defined section over
$T^*(M/B)$.

\emph{Replacing $\epsilon$ by $2\epsilon$.} Let $\tilde A(\epsilon)=A(2\epsilon)$,
$0<\epsilon\le\frac12$. Its kernel is $\epsilon^{-n}\tilde K_2(\epsilon,z,(z-z')/
\epsilon)$ with $\tilde K_2(\epsilon,z,W)=2^{-n}K(2\epsilon,z,W/2)$, so $\tilde A$ is
again semiclassical, and
\begin{equation*}
  \sigsl(\tilde A)(z,\zeta)=2^{-n}\int e^{-iW\cdot\zeta}K(0,z,W/2)\,dW=\hat K(0,z,2\zeta)
  =\sigsl(A)(z,2\zeta).
\end{equation*}
The oscillatory characterization shows the same thing:
$e^{-if/\epsilon}A(2\epsilon)e^{if/\epsilon}=e^{-i(2f)/(2\epsilon)}A(2\epsilon)
e^{i(2f)/(2\epsilon)}\to\sigsl(A)(z,2df)$. So the same family of operators has a
different symbol when the parameter is changed: the symbol is composed with the
fibre dilation by $2$. This is the case $T\equiv2$ of \cref{scindex:rem:slT}. The
invariant object is the covector $d(f/\epsilon)$ of the phase, which explains why
$\sigsl$ lives naturally on ${}^{\mathrm{sl}}T^*(M/B)$. The spaces
$\Psisl{-\infty}$, $\epsilon\Psisl{-\infty}$ and $\Psiinf$ are unchanged by the
reparametrization, and since the dilations by $s\in[1,2]$ are homotopic to the
identity through bundle isomorphisms, no K-theoretic statement of the chapter
depends on the choice.

%----------------------------------------------------------------------
\subsection*{\Cref{scindex:ex:ground-state}}

\emph{Idea:} $\Gamma$ is an isometry, so $\Gamma(\cdot)\Gamma^*$ is multiplicative;
and $\Gamma$ is one more isometric embedding to which \cref{scindex:lem:class}(i)
applies.

Let $k$ be the rank of $\mathcal V$, $\nu$ the density on $Z$, and give $\mathcal V$
the density $\varpi^*\nu\otimes dv$, $dv$ the Lebesgue measure of the metric on the
fibres. The fibrewise normalized Gaussian is $\gamma(v)=\pi^{-k/4}e^{-\abs v^2/2}$,
so $\int_{\mathcal V_z}\gamma^2dv=1$. In a local orthonormal frame,
$\mathcal V|_U\cong U\times\R^k$ and $\gamma(z,w)=\pi^{-k/4}e^{-\abs w^2/2}$ does
not depend on $z$.

\emph{$\Gamma$ is an isometry.} For $u\in L^2(Z;E)$,
$\norm{\Gamma u}^2=\int_Z\abs{u(z)}^2\int_{\mathcal V_z}\gamma^2\,dv\,\nu=\norm u^2$.
Its adjoint is $(\Gamma^*w)(z)=\int_{\mathcal V_z}\gamma(v)w(v)\,dv$, and
$\Gamma^*\Gamma=\Id$.

\emph{The embedding.} For $A\in\Psido{-\infty}(Z;E)$ with kernel $A(z,z')$, the
kernel of $\Gamma A\Gamma^*$ is $\gamma(v)A(\varpi v,\varpi v')\gamma(v')$ (with
respect to the density on $\mathcal V$). In local frames it is
$\pi^{-k/2}e^{-(\abs w^2+\abs{w'}^2)/2}A(z,z')$: smooth, and Schwartz in the fibre
variables $(w,w')$ uniformly with all derivatives, i.e.\ in
$\Psido{-\infty}_{\mathcal S}(\mathcal V;\varpi^*E)$ (\cref{scindex:rem:schwartz}
with $M=Z$ and $B$ a point). The map $A\mapsto\Gamma A\Gamma^*$ is linear,
\begin{equation*}
  (\Gamma A\Gamma^*)(\Gamma B\Gamma^*)=\Gamma AB\Gamma^*,\qquad
  (\Gamma A\Gamma^*)^*=\Gamma A^*\Gamma^*,
\end{equation*}
by $\Gamma^*\Gamma=\Id$, and injective since $\Gamma^*(\Gamma A\Gamma^*)\Gamma=A$.
So it embeds $\Psido{-\infty}(Z;E)$ as a $*$-subalgebra (without unit) of the
operators with Schwartz kernels.

\emph{Invertible elements.} If $g=\Id+A\in\Ginf(Z;E)$ with $g^{-1}=\Id+A'$, put
$\hat g=\Id+\Gamma A\Gamma^*=\Gamma g\Gamma^*+(\Id-\Gamma\Gamma^*)$. Then
$\hat g\,\widehat{g^{-1}}=\Id+\Gamma(A+A'+AA')\Gamma^*=\Id$ and similarly in the
other order, so $\hat g\in\Id+\Psido{-\infty}_{\mathcal S}$ is invertible, and
$g\mapsto\hat g$ is a group homomorphism.

\emph{Classes.} Since $K^{-1}(\mathrm{pt})=[\mathrm{pt};\Ginf]=0$ (every element of
$\Ginf$ is joined to one in a connected group $\GL(V)$ by \cref{bott:lem:stable}),
the statement for a point is empty; the content is the statement with parameters,
which we prove. Let $X$ be a compact manifold and $g:X\to\Ginf(Z;E)$ smooth, so
$[g]\in K^{-1}(X)$ by \ref{scindex:K-Ginf} (if $\dim Z=0$, $\Ginf(Z;E)$ is a group $\GL(N,\C)$ and $[g]$ is its stable class; the argument below is unchanged). Then
$\hat g$ is an element of $\Id+\Psido{-\infty}_{\mathcal S}$ for the fibration
$\mathcal V\times X\to X$, invertible on every fibre, and its class is defined by
\cref{scindex:def:class} through an isometric embedding $J:L^2(\mathcal V;\varpi^*E)\to
L^2(\Torus^N;\C^R)$ as in \cref{scindex:rem:schwartz}, independent of $x\in X$:
\begin{equation*}
  \rho_J(\hat g)=\Id+J\Gamma A\Gamma^*J^*=IgI^*+\Id-II^*,\qquad I=J\Gamma .
\end{equation*}
$I$ is an isometry, and $T\mapsto ITI^*=J(\Gamma T\Gamma^*)J^*$ maps
$\Psido{-\infty}(Z;E)$ continuously into $\Psido{-\infty}(\Torus^N;\C^R)$, by the
first part and the admissibility of $J$. Moreover $I$ maps $\CI(Z;E)$ into
$\CI(\Torus^N;\C^R)$: if $u\in\CI(Z;E)$, $u\ne0$, the rank one operator
$u\otimes u^*$ is smoothing, so $I(u\otimes u^*)I^*=(Iu)\otimes(Iu)^*$ has a smooth
kernel $(Iu)(y)\overline{(Iu)(y')}$; choosing $y_0$ with $(Iu)(y_0)\ne0$ (as a
Lebesgue point), $Iu(y)=k(y,y_0)/\overline{(Iu)(y_0)}$ is smooth. So
\cref{scindex:lem:class}(i) applies: $g\mapsto IgI^*+\Id-II^*$ induces the
canonical identification $[X;\Ginf(Z;E)]\cong[X;\Ginf(\Torus^N;\C^R)]$ of
\ref{scindex:K-Ginf}. Hence $[\hat g]=[\rho_J(\hat g)]=[g]$ in $K^{-1}(X)$: the
embedding $A\mapsto\Gamma A\Gamma^*$ preserves the classes.

%----------------------------------------------------------------------
\subsection*{\Cref{scindex:ex:circle}}

\emph{Idea:} choose the cut-off in $\xi$ so that at $\epsilon=1$ it is exactly the
projection $\Pi_>$; then the quantized Milnor idempotent at $\epsilon=1$ is $p_0$
plus a rank one idempotent, and the deformation in $\epsilon$ does not change the
relative index.

\emph{The symbol.} Choose $\chi\in\CI(\R)$, non-decreasing, with $\chi=0$ on
$(-\infty,0]$ and $\chi=1$ on $[1,\infty)$, and put
\begin{equation*}
  a(\theta,\xi)=1-\chi(\xi)+\chi(\xi)e^{-i\theta}.
\end{equation*}
Then $a=1$ for $\xi\le0$ (in particular for $\xi<-1$) and $a=e^{-i\theta}$ for
$\xi\ge1$, so $a$ is homogeneous of degree $0$ and invertible for $\abs\xi\ge1$,
equal there to $\sigma_0(P_1)$ by \eqref{families:eq:P1}. As in the proof of
\cref{scindex:thm:analytic=sl}, $x=[(\C,\C,a)]\in\Kc(T^*\Torus)$ is the class of
$\sigma_0(P_1)$. Put $b(\theta,\xi)=1-\chi(\xi)+\chi(\xi)e^{i\theta}$. Then
\begin{equation*}
  ba=ab=(1-\chi)^2+\chi^2+2\chi(1-\chi)\cos\theta,\qquad
  r_+=r_-=r=2\chi(1-\chi)(1-\cos\theta),
\end{equation*}
which is supported in $\{0\le\xi\le1\}$; so $b$ satisfies \eqref{scindex:eq:b}, and
on $E_-\oplus E_+=\C\oplus\C$, by \eqref{scindex:eq:milnor},
\begin{equation*}
  e_{a,b}=\begin{pmatrix}1-r^2&(1+r)ar\\rb&r^2\end{pmatrix},\qquad
  p_0=\begin{pmatrix}1&0\\0&0\end{pmatrix}.
\end{equation*}

\emph{A class of semiclassical operators on the circle.} Write $e_k=e^{ik\theta}$,
$D=D_\theta$, and for a bounded smooth $g$ on $[0,1]\times\R$ let $g(\epsilon,\epsilon D)$
be the Fourier multiplier $e_k\mapsto g(\epsilon,\epsilon k)e_k$. Let $\mathcal T$ be
the space of finite sums $\sum_je^{ij\theta}g_j(\epsilon,\epsilon D)$ with
$g_j\in\CI([0,1]\times\R)$ compactly supported in $\xi$. By Poisson summation, the
kernel of $g(\epsilon,\epsilon D)$ is
\begin{equation*}
  \frac1{2\pi}\sum_kg(\epsilon,\epsilon k)e^{ik(\theta-\theta')}
  =\sum_{m\in\Z}\epsilon^{-1}\check g\Bigl(\epsilon,\frac{\theta-\theta'-2\pi m}\epsilon\Bigr),
  \qquad\check g(\epsilon,W)=\frac1{2\pi}\int e^{iW\xi}g(\epsilon,\xi)\,d\xi,
\end{equation*}
with $\check g$ Schwartz in $W$ uniformly with all derivatives. For
$\theta-\theta'\in(-\pi,\pi]$ the terms with $m\ne0$, and the term $m=0$ away from
the diagonal, are $O(\epsilon^\infty)$ with all derivatives. Hence
$g(\epsilon,\epsilon D)\in\Psisl{-\infty}(\Torus)$ with $\sigsl=g(0,\xi)$, and by
\ref{scindex:S-symbol} (multiplication by $e^{ij\theta}$) every element of
$\mathcal T$ is in $\Psisl{-\infty}$ with $\sigsl\bigl(\sum e^{ij\theta}g_j\bigr)=
\sum_je^{ij\theta}g_j(0,\xi)$. Moreover, if $h$ is smooth on $[0,1]\times\R$ with
bounded derivatives, then
\begin{equation}\label{solsemiclassical-index:eq:shift}
  h(\epsilon,\epsilon D)\,e^{il\theta}=e^{il\theta}h(\epsilon,\epsilon D+\epsilon l),
\end{equation}
so products $e^{ij\theta}h(\epsilon,\epsilon D)\cdot e^{il\theta}g(\epsilon,\epsilon D)$
and $e^{il\theta}g(\epsilon,\epsilon D)\cdot e^{ij\theta}h(\epsilon,\epsilon D)$ with
$g$ compactly supported in $\xi$ are again in $\mathcal T$, with symbol at $\epsilon=0$
the product of the symbols. For $\epsilon\in[\epsilon_0,1]$ only the finitely many
modes $\abs k\le C/\epsilon_0$ contribute to an element of $\mathcal T$, so its kernel
is a finite trigonometric sum with coefficients smooth in $\epsilon$.

\emph{The quantized pair.} Let
\begin{equation*}
  A(\epsilon)=1-\chi(\epsilon D)+e^{-i\theta}\chi(\epsilon D),\qquad
  B(\epsilon)=1-\chi(\epsilon D)+e^{i\theta}\chi(\epsilon D).
\end{equation*}
By \eqref{solsemiclassical-index:eq:shift}, with $\chi=\chi(\epsilon D)$ and
$\chi_\pm=\chi(\epsilon D\pm\epsilon)$,
\begin{align*}
  R_+&=1-BA=1-(1-\chi)^2-\chi_-\chi-e^{i\theta}\chi(1-\chi)-e^{-i\theta}(1-\chi_-)\chi,\\
  R_-&=1-AB=1-(1-\chi)^2-\chi_+\chi-e^{-i\theta}\chi(1-\chi)-e^{i\theta}(1-\chi_+)\chi .
\end{align*}
All multipliers here are compactly supported in $\xi$ for $\epsilon\in[0,1]$
(they vanish for $\xi\le0$, where $\chi=0$, and for $\xi\ge2$, where
$\chi=\chi_\pm=1$), so $R_\pm\in\mathcal T$, with $\sigsl(R_\pm)=1-(1-\chi)^2-\chi^2-
2\chi(1-\chi)\cos\theta=r$. The entries $-R_-^2$, $(1+R_-)AR_+$, $R_+B$, $R_+^2$ of
$e(\epsilon)-p_0$, $e(\epsilon)=e_{A(\epsilon),B(\epsilon)}$, are therefore in
$\mathcal T$, with semiclassical symbols the entries $-r^2$, $(1+r)ar$, $rb$, $r^2$
of $e_{a,b}-p_0$. By \cref{scindex:lem:milnor}, $e(\epsilon)$ is an idempotent for
every $\epsilon\in(0,1]$. So $e(\epsilon)$ is a semiclassical quantization of
$e_{a,b}$ in the sense of \cref{scindex:prop:quantize}, and by
\cref{scindex:def:Indsl},
\begin{equation*}
  \Indsl{}(x)=d\bigl(e(\epsilon),p_0\bigr)\quad\text{for small }\epsilon>0.
\end{equation*}

\emph{Deformation to $\epsilon=1$.} For $\epsilon\in[\epsilon_0,1]$ the family
$e(\epsilon)-p_0$ has smooth kernels depending smoothly on $\epsilon$, so it is a
smooth family of compactly supported idempotent perturbations of $p_0$ and, by
\ref{scindex:d-homotopy}, $d(e(\epsilon),p_0)=d(e(1),p_0)$.

\emph{The case $\epsilon=1$.} Since $\chi(k)=0$ for $k\le0$ and $\chi(k)=1$ for
$k\ge1$, $\chi(D)=\Pi_>$, and $A(1)=P_1$, $B(1)=1-\Pi_>+e^{i\theta}\Pi_>$. On the
basis: $A(1)e_k=e_k$ for $k\le0$ and $A(1)e_k=e_{k-1}$ for $k\ge1$; $B(1)e_k=e_k$ for
$k\le0$ and $B(1)e_k=e_{k+1}$ for $k\ge1$. Hence $A(1)B(1)=\Id$, and
$B(1)A(1)e_k=e_k$ except $B(1)A(1)e_1=B(1)e_0=e_0$. Thus
\begin{equation*}
  R_-=0,\qquad R_+=\frac1{2\pi}(e_1-e_0)\otimes\bar e_1:\ u\mapsto
  \frac1{2\pi}\inner u{e_1}(e_1-e_0),
\end{equation*}
a rank one operator with $R_+^2=R_+$ (as $\inner{e_1-e_0}{e_1}=2\pi$). Moreover
$AR_+=R_-A=0$ and $R_+B=BR_-=0$. By \eqref{scindex:eq:milnor},
\begin{equation*}
  e(1)=\begin{pmatrix}1&0\\0&R_+\end{pmatrix}=p_0+\Pi,\qquad
  \Pi=\begin{pmatrix}0&0\\0&R_+\end{pmatrix},\quad p_0\Pi=\Pi p_0=0,
\end{equation*}
and $\Pi$ is a smoothing idempotent of rank one. By \ref{scindex:d-finite},
$d(e(1),p_0)=[\Ran\Pi]=[\C]=1$. Therefore
\begin{equation*}
  \Indsl{}\bigl([\sigma_0(P_1)]\bigr)=1,
\end{equation*}
in agreement with $\ind(P_1)=1$ and \cref{scindex:thm:analytic=sl}. (Alternatively,
by \cref{solsemiclassical-index:lem:d-trace} and the computation of
\cref{scindex:ex:milnor-trace}, $d(e(\epsilon),p_0)=\Tr R_+-\Tr R_-=\sum_k
\bigl(\chi(\epsilon k)\chi(\epsilon k+\epsilon)-\chi(\epsilon k-\epsilon)\chi(\epsilon k)\bigr)$,
a telescoping sum equal to $1$ for every $\epsilon$.)
